% example.tex —— 演示：汉化后的 biblatex verbose-trad2（脚注中的 ibid./idem/op. cit. 均为中文）
% 编译流程：xelatex example -> biber example -> xelatex example -> xelatex example

\begin{filecontents*}[overwrite]{refs.bib}
@book{zhang2010,
	author    = {张伟},
	title     = {民法总论},
	publisher = {法律出版社},
	address   = {北京},
	year      = {2010},
}

@book{zhang2015,
	author    = {张伟},
	title     = {合同法研究},
	publisher = {中国人民大学出版社},
	address   = {北京},
	year      = {2015},
}

@book{li2020,
	author     = {李娜},
	title      = {比较法导论},
	translator = {王芳},
	publisher  = {北京大学出版社},
	address    = {北京},
	year       = {2020},
}

@article{wang2019,
	author    = {王刚},
	title     = {论正当防卫},
	journal   = {法学研究},
	volume    = {3},
	number    = {2},
	year      = {2019},
	pages     = {10--25},
}

@collection{chen2018,
	editor    = {陈华},
	title     = {宪法学文集},
	publisher = {法律出版社},
	address   = {北京},
	year      = {2018},
}

@phdthesis{zhao2021,
	author = {赵磊},
	title  = {明代财政制度研究},
	school = {清华大学},
	address = {北京},
	year   = {2021},
}

@book{smith2001,
	author    = {Smith, John},
	title     = {The History of Things},
	publisher = {Oxford University Press},
	address   = {Oxford},
	year      = {2001},
}

@book{liu2010,
	author    = {刘强},
	title     = {宪法学原理},
	publisher = {中国政法大学出版社},
	address   = {北京},
	year      = {2010},
}

@book{liu2018,
	author    = {刘强},
	title     = {行政法研究},
	publisher = {法律出版社},
	address   = {北京},
	year      = {2018},
}
\end{filecontents*}

\documentclass{ctexart}

\usepackage[style=chinese-verbose-trad2,
            language=chinese,
            ibidpage=true,
            backend=biber]{biblatex}

\addbibresource{refs.bib}

\begin{document}

\section*{汉化 verbose-trad2 脚注示例}

下面一组脚注依次演示了“完整引用—同作者异书(idem)—同上(ibidem)—前引书(op.~cit.)”
的完整链条，所有提示词均已被汉化。\footcite[25]{zhang2010}

随后引用同一作者的另一部作品（首次出现，给完整信息）。%
\footcite{zhang2015}

相邻脚注再次引用同一文献，输出“同上”（ibidem）。\footcite[30]{zhang2015}

再引用一位不同作者的译著，会给出完整信息（含“译”）。\footcite{li2020}

隔了几条之后再次引用张伟的《民法总论》，作者相同但非紧邻，
提示词变为“前引书”（op.~cit.）。\footcite[12]{zhang2010}

期刊论文的完整脚注格式如下。\footcite{wang2019}

编著类文献的责任方式显示为“编”。\footcite{chen2018}

学位论文也能正确本地化。\footcite{zhao2021}

西文文献同样适用该样式，仅在术语处切换为中文提示词。\footcite[15]{smith2001}

再次引用前文期刊论文时使用“前引书”。\footcite[20]{wang2019}

行内使用 \verb|\autocite| 也会落入脚注：\autocite{zhang2010}

先把刘强《宪法学原理》作为完整引用\footcite{liu2010}；随后在同一脚注内先引其
《行政法研究》、再引《宪法学原理》，后者因同作者而缩写为“同上”（idem）：%
\footnote{如\footcite{liu2018}\footcite{liu2010}。}

\bigskip
\noindent 说明：以上脚注中的“同上”对应 \texttt{ibidem}，“前引书”对应
\texttt{op.~cit.}；同作者异书情形由 \texttt{idem} 控制。这些字符串均定义在
随附的 \texttt{chinese.lbx} 中，可按需修改。

\printbibliography

\end{document}
